% GENERATED_BY: oc_core_monograph_translation_draft_pipeline_v1.py
% AI_ASSISTED_TRANSLATION_DRAFT: TRUE
% THEOREM_REVIEW_STATUS: PENDING
% THEOREM_REVIEW_MODE: FORMAL_AI_GATE__CODEX_CLI_CHATGPT
% THEOREM_REVIEW_REVIEWER_ID: 
% THEOREM_REVIEW_AUTHORITY_CLASS: 
% THEOREM_REVIEWED_AT: 
% THEOREM_REVIEW_DOSSIER_REF: 
% SOURCE_LANGUAGE: EN
% TARGET_LANGUAGE: RU
% SOURCE_EN_REF: content/19_oc_core_1_3_foundational_consistency.tex
\section{Core 1.3 Foundational Consistency Dossier}
\label{sec:oc13-foundational-consistency}

Core 1.3 adds an explicit consistency layer to the inherited monograph shell.
Its purpose is not to replace the formal chapters with editorial commentary,
but to make three things readable in one place:

\begin{itemize}
\item why the present hierarchy is fixed as \(K_0\) through \(K_{12}\) rather
      than by an arbitrary cutoff;
\item which statements in the volume are axiomatic, derived, structural, or
      interpretive;
\item how the master monograph and the narrower journal-core extraction are
      related without silent drift in scope or notation.
\end{itemize}

\subsection{Why the hierarchy is fixed as \texorpdfstring{$K_0$--$K_{12}$}{K0-K12}}

Core 1.2 already contains enough structure to force a hierarchy larger than
\(K_{10}\). The point is not that higher labels are rhetorically attractive.
The point is that the исходным корпусом itself already distinguishes:

\begin{enumerate}
\item lower-level continua whose axes, thresholds, and flows are internal to a
      given model class;
\item theoretical and meta-theoretical continua that organise theories and
      model-spaces themselves;
\item a further stratum in which operators, meta-logics, and model-spaces
      become explicit objects of structural evolution;
\item and an upper bounding layer in which global coherence conditions across
      such higher-order structures must themselves be represented.
\end{enumerate}

That is why Core 1.3 fixes the hierarchy as \(K_0,\dots,K_{12}\) rather than
as \(K_0,\dots,K_{10}\). Stopping at \(K_{10}\) leaves the публичный исходным корпусом
internally inconsistent, because later modules already require
\(K_{11}\) and \(K_{12}\) to state the cross-level transitions and
upper-boundary conditions that the framework claims to discuss.

\subsection{Scientific role of the upper levels}

\paragraph{Level \texorpdfstring{$K_{11}$}{K11}.}
\(K_{11}\) is the first level at which formal ontologies, operator systems,
meta-logics, and model-spaces are themselves treated as structured continua
subject to admissibility, transformation, and collapse conditions.
This level is needed once one studies not only theories, but the evolution of
the spaces in which theories and theory-building operators live.

\paragraph{Level \texorpdfstring{$K_{12}$}{K12}.}
\(K_{12}\) is not introduced as a mystical completion layer.
It is the minimally specified upper coherence stratum required when one needs
to discuss global compatibility, boundedness, and failure modes across families
of \(K_{11}\)-objects.
Its role is structural rather than triumphalist: it prevents the hierarchy from
pretending that there is no formal problem above recursively evolving
meta-theoretical spaces.

\subsection{Statement classes}

To keep the monograph readable without blurring proof status, Core 1.3 uses the
following statement discipline throughout the release:

\begin{itemize}
\item \textbf{Axiomatic statements} fix primitive assumptions, admissibility
      conditions, or identity rules.
\item \textbf{Derived statements} are теорему-bearing rows with explicit support
      dependence and proof obligations.
\item \textbf{Structural statements} describe the organisation of the framework,
      the hierarchy, and the mapping between modules or levels.
\item \textbf{Interpretive statements} explain what the formal structure means,
      which reader it is for, and what later projections may test.
\end{itemize}

The journal-core extraction is allowed to narrow this body, but not to blur
these classes into one another.
An interpretive sentence may explain a теорему; it may not replace the теорему.
A structural remark may motivate a projection; it may not silently serve as
proof of that projection.

\subsection{Proof-bearing and не-утверждение layers}

The master monograph therefore separates three layers explicitly:

\begin{enumerate}
\item the inherited formal backbone that still bears the structure of the core
      theory;
\item the Core 1.3 repair layer that clarifies witness discipline, теорему
      fate, chronology, and публичныйation boundary;
\item the bounded не-утверждение layer that states what the present release does not
      yet defend, even when the broader research program still pursues it.
\end{enumerate}

This separation matters for hostile review.
A reviewer should be able to ask, for any sentence that sounds strong:
is it axiomatic, derived, structural, or interpretive?
If it is derived, what supports it?
If it is structural or interpretive, what exactly is it not claiming?

\subsection{Relation between master monograph and journal core}

The master monograph is the canonical scientific source because it carries the
full theory, the reader ladder, the appendices, and the explicit публичныйation
boundary in one place.
The journal-core статья is narrower by design.
Its role is to extract one defensible теорему-bearing route for editorial
inspection without pretending that the whole monograph is already compressed
into that shorter artifact.

For that reason the extraction rule is one-way:

\[
\text{Master Monograph} \longrightarrow \text{Journal Core},
\]

with explicit preservation of notation, не-утверждение boundaries, and теорему scope.
The shorter статья may omit detail.
It may not enlarge the claim.

\subsection{Reader-facing consequence}

This chapter exists because scientific rigor and readability do not pull in
opposite directions here.
The more explicit the dossier becomes about hierarchy, statement class, and
proof boundary, the easier it becomes for three readers at once:

\begin{itemize}
\item the editor who needs a clean overview of what is actually defended;
\item the formal reader who needs exact scope and dependency discipline;
\item the non-specialist scientific reader who needs a coherent narrative entry
      point before diving into the full technical shell.
\end{itemize}

\clearpage
